% ============================================================
%  Chương 5 – Tính năng nâng cao
%             Làn sóng thương mại điện tử thứ tư  (1 điểm)
%  Trạng thái: ĐÃ HOÀN THIỆN
% ============================================================

\section*{Giới thiệu -- Làn sóng thương mại điện tử thứ tư}

UniJob không chỉ là một bảng thông báo việc làm. Nhóm thiết kế và triển khai 5 tính năng
thể hiện rõ vị trí của nền tảng trong làn sóng thứ tư.

\section{AI-Powered Job Recommendation}

\subsection*{Mô tả tính năng}

Section ``Đề xuất cho bạn'' trên trang chủ hiển thị tối đa 6 công việc được cá nhân hóa
cho mỗi sinh viên sau khi đăng nhập. Giao diện gồm lưới 3 cột với badge lý do gợi ý.

\subsection*{Cơ chế kỹ thuật}

Thuật toán client-side trong \texttt{suggest.service.ts} chấm điểm từng công việc
(thang 0--105 điểm):

\begin{table}[H]
\centering
\caption{Bảng trọng số thuật toán AI Matching}
\begin{tabularx}{\textwidth}{|p{4cm}|c|X|}
\hline
\textbf{Yếu tố} & \textbf{Điểm} & \textbf{Logic} \\
\hline
Faculty match    & 40 pts & \texttt{job.faculty === user.faculty} \\
\hline
Skills/tags overlap & 30 pts & tỷ lệ giao tag kỹ năng \\
\hline
Category affinity & 20 pts & người dùng từng ứng tuyển category này \\
\hline
Recency          & 10 pts & decay tuyến tính trong 7 ngày \\
\hline
Urgent bonus     & +5 pts & job có flag \texttt{isUrgent = true} \\
\hline
\end{tabularx}
\end{table}

Loại trừ: công việc của chính user, đã ứng tuyển, không ở trạng thái \texttt{open}.

\subsection*{Tính năng}

Thuật toán gợi ý cá nhân hóa là đặc trưng của \textit{AI-driven commerce}: nền tảng không
chỉ hiển thị danh sách, mà hiểu ngữ cảnh người dùng (khoa, kỹ năng, lịch sử) để đề xuất
chủ động. Đây là sự khác biệt căn bản so với marketplace thế hệ trước dựa trên keyword search.

\section{Digital Certification -- CV Passport}

\subsection*{Mô tả tính năng}

Sinh viên có thể xuất \textbf{Chứng nhận Kinh nghiệm Thực tế} dưới dạng file PDF từ
trang \texttt{/cv-export}. Chứng nhận bao gồm: tên, khoa, MSSV; 3 stat box
(số công việc / điểm đánh giá / tổng giờ làm); danh sách dự án đã hoàn thành.

\subsection*{Cơ chế kỹ thuật}

Sử dụng \textbf{Canvas 2D API} để vẽ chứng nhận trực tiếp trên browser (không qua server),
sau đó nhúng vào jsPDF:

\begin{enumerate}
  \item \texttt{generateCVCertificate()} tạo canvas \texttt{794×1123px} (A4 96dpi).
  \item Scale 3× để đảm bảo độ phân giải sắc nét.
  \item Vẽ header xanh, logo, tên user, stat boxes, danh sách job từ Firestore.
  \item Export sang PNG $\rightarrow$ nhúng vào jsPDF $\rightarrow$ \texttt{pdf.save()}.
\end{enumerate}

Placeholder QR code dự kiến tích hợp trong sprint tiếp theo để liên kết đến
URL xác minh trực tuyến.

\subsection*{Tính năng}

CV Passport thuộc xu hướng \textit{Portable Reputation \& Digital Credentials} -- nơi
danh tiếng không nằm trong CV Word mà được gắn trực tiếp vào hành vi giao dịch thực tế.
Trong nền kinh tế gig economy, chứng nhận minh bạch này có giá trị hơn hồ sơ tự viết.

\section{Gamification \& Trust Score -- Hệ thống Điểm Uy Tín}

\subsection*{Mô tả tính năng}

Mỗi user có \textbf{Điểm Uy Tín} (\texttt{ratingScore}) hiển thị trên profile và job card.
Điểm được cập nhật sau mỗi lần đánh giá 1--5 sao khi hoàn thành task.

\subsection*{Cơ chế kỹ thuật}

\begin{itemize}
  \item Collection \texttt{/ratings/\{ratingId\}} lưu mỗi đánh giá: điểm, comment,
    fromUserId, toUserId.
  \item Security Rule ngăn self-rating (\texttt{fromUserId !== toUserId}).
  \item \texttt{ratingScore} trên user được tính trung bình cộng dần (incremental average).
  \item Lịch sử hoạt động và \texttt{workHistory} tích lũy theo thời gian.
\end{itemize}

\subsection*{Tính năng}

Hệ thống reputation là nền tảng của \textit{Trust Economy} -- mô hình kinh doanh mà giao
dịch xảy ra dựa trên danh tiếng tích lũy thay vì hợp đồng cứng nhắc. Airbnb, Uber và
Fiverr đều xây dựng trên nguyên lý này.

\section{Smart Matching for Collaboration}

\subsection*{Mô tả tính năng}

Khi poster đăng task mới, hệ thống có thể gợi ý \textbf{sinh viên đã từng hợp tác tốt}
dựa trên lịch sử \texttt{workHistory} -- ưu tiên người có điểm đánh giá cao và đã làm
cùng category trước đây.

\subsection*{Cơ chế kỹ thuật}

\begin{itemize}
  \item Collection \texttt{/workHistory/} lưu \{userId, jobId, jobTitle, category, completedAt\}.
  \item Khi tạo job mới, query lịch sử để tìm worker phù hợp nhất (cùng category,
    ratingScore $\geq$ 4).
  \item Tính năng đang ở giai đoạn thiết kế, dự kiến hoàn thành Sprint 4.
\end{itemize}

\subsection*{Tính năng}

\textit{Relationship-based commerce} -- khai thác mạng lưới quan hệ đã được xác minh
thay vì tìm kiếm người lạ. Đây là xu hướng của LinkedIn Gigs, Contra và các nền tảng
freelance thế hệ mới.

\section{Real-time Campus Marketplace}

\subsection*{Mô tả tính năng}

\begin{itemize}
  \item \textbf{Urgent Mode:} Khi đăng task có flag \texttt{isUrgent = true}, task
    hiển thị badge ``Khẩn cấp'' màu cam nổi bật và được AI Matching ưu tiên +5 điểm.
  \item \textbf{Real-time state sync:} Zustand store + Firestore \texttt{onSnapshot}
    đảm bảo trạng thái task (open/assigned/completed) cập nhật ngay lập tức cho tất cả
    người dùng đang xem.
  \item \textbf{Toast notifications:} Mọi hành động (ứng tuyển, chấp nhận, hoàn thành)
    đều phản hồi tức thì qua toast -- không cần reload trang.
\end{itemize}

\subsection*{Tính năng}

\textit{Real-time \& Hyper-local Commerce} -- mô hình mà giá trị được tạo ra từ sự phù
hợp cung--cầu \textbf{ngay lập tức trong một vùng địa lý hẹp}. Đây là bản chất của Grab,
Gojek hay các nền tảng delivery: không gian campus chính là ``hyper-local'' nhỏ nhất và
hiệu quả nhất để kiểm chứng mô hình.

\section*{Tổng kết}

\begin{table}[H]
\centering
\caption{Bảng tổng hợp 5 tính năng nâng cao và vị trí trong làn sóng thứ tư}
\renewcommand{\arraystretch}{1.2}
\begin{tabularx}{\textwidth}{|X|p{4cm}|p{3.5cm}|}
\hline
\textbf{Tính năng} & \textbf{Trụ cột 4th Wave} & \textbf{Trạng thái} \\
\hline
AI Job Recommendation  & AI Personalization      & Đã triển khai \\
\hline
CV Passport (PDF)      & Digital Credentials     & Đã triển khai \\
\hline
Gamification + Trust Score & Trust Economy      & Đã triển khai \\
\hline
Smart Matching         & Relationship Commerce   & Đang phát triển \\
\hline
Real-time Marketplace  & Hyper-local / Real-time & Đã triển khai \\
\hline

\end{tabularx}
\end{table}

Các tính năng trên không phải bổ sung trang trí -- chúng là \textbf{điều kiện cần} để một
campus marketplace vận hành được trong thời đại sinh viên kỳ vọng trải nghiệm cá nhân hóa,
minh bạch và tức thời như các ứng dụng tiêu dùng họ dùng hàng ngày.
